\chapter{Flat continuity and exceptional slicing levels}\label{sec:slice-flat-continuity}
The lower and upper slices are meaningful as distributional currents even at exceptional levels. Their flat norm is controlled by the mass in a thin slab. Throughout this section $T$ is a compactly supported normal $m$-current in $\mathbb R^N$, $m\geq1$, and $f$ is locally Lipschitz near its support. Flat norms are the real-current norms of Definition \ref{def:real-current-flat-norm}. These deductions refine the distinction between the two one-sided slices in \cite{Sim18}, equations (4.6)--(4.10), pp.~184--185.

\begin{definition}[Lower and upper distributional slices]\label{def:two-sided-current-slices}
For every real $t$, set
\[
 B_t^- =\partial(T\llcorner\{f<t\})-(\partial T)\llcorner\{f<t\},
 \qquad
 B_t^+ =-\partial(T\llcorner\{f>t\})+(\partial T)\llcorner\{f>t\}.
\]
Since $\partial T=\partial(T\llcorner\mathbb R^N)$, the upper slice can also be written
\[
 B_t^+=\partial(T\llcorner\{f\leq t\})-(\partial T)\llcorner\{f\leq t\}.
\]
Neither formula alone guarantees finite mass at every level. They agree with the usual geometric slice at almost every level when $T$ is integral. The signs agree with \cite{Sim18}, equations (4.6)--(4.7).
\end{definition}

\begin{proposition}[Support and finite flat norm at every level]\label{prop:all-level-slice-support-flat}
Both $B_t^-$ and $B_t^+$ are compactly supported currents with
\[
 \operatorname{spt}B_t^\pm\subset\operatorname{spt}T\cap\{f=t\},
 \qquad
 \mathcal F_{\mathbb R}(B_t^\pm)
 \leq\mathbf M(T)+\mathbf M(\partial T).
\]
Finite flat norm here does not imply finite mass. Compare the support assertion in \cite{Sim18}, equation (4.8).
\end{proposition}
\begin{proof}
For the lower slice use the flat decomposition with filling $T\llcorner\{f<t\}$ and residual $-(\partial T)\llcorner\{f<t\}$. Both have finite mass and compact support in $\operatorname{spt}T$, since $\operatorname{spt}\partial T\subset\operatorname{spt}T$. This proves the norm bound. Use the sublevel formula with $\leq$ for the upper slice.

On the open set $\{f<t\}$, the lower slice acts as $\partial T-\partial T=0$; on $\{f>t\}$, both terms vanish. These statements are local: every test compactly supported in either open set has that same property on the support of its derivative. A partition of unity for a compact test support disjoint from the level extends this to all tests off $\{f=t\}$. The same argument applies to the upper slice. Tests off $\operatorname{spt}T$ vanish as well, proving the inclusion.
\end{proof}

\begin{proposition}[A slab controls the flat change of a slice]\label{prop:slice-flat-slab-estimate}
For $s<t$,
\[
 B_t^- -B_s^-
 =\partial(T\llcorner\{s\leq f<t\})
       -(\partial T)\llcorner\{s\leq f<t\},
\]
and hence, with $\mu=\|T\|+\|\partial T\|$,
\[
 \mathcal F_{\mathbb R}(B_t^- -B_s^-)
 \leq\mu(\{s\leq f<t\}).
\]
The analogous upper-slice bound uses $\{s<f\leq t\}$. This is a flat-norm consequence of the restriction identities in \cite{Sim18}, equations (4.6)--(4.7).
\end{proposition}
\begin{proof}
Subtract the two lower-sublevel definitions and use the exact indicator identity $\chi_{\{f<t\}}-\chi_{\{f<s\}}=\chi_{\{s\leq f<t\}}$. The two resulting terms are an admissible real flat decomposition. Their mass sum is at most the displayed slab measure. Subtracting the closed-sublevel formulas for the upper slices gives the stated alternative. The choice of strict and nonstrict endpoints matters at levels carrying mass.
\end{proof}

\begin{theorem}[One-sided continuity and a countable jump set]\label{thm:slice-flat-one-sided}
The map $t\mapsto B_t^-$ is left continuous in real flat norm and has right limit $B_t^+$ at every $t$. The map $t\mapsto B_t^+$ is right continuous and has left limit $B_t^-$. Their difference is
\[
 B_t^+-B_t^-=
 \partial(T\llcorner\{f=t\})-(\partial T)\llcorner\{f=t\}.
\]
Both maps are continuous at every $t$ such that $\mu(\{f=t\})=0$, where $\mu=\|T\|+\|\partial T\|$. Consequently any discontinuities lie in a countable set. These conclusions explain the one-sided conventions in \cite{Sim18}, p.~184.
\end{theorem}
\begin{proof}
If $s\uparrow t$, the sets $\{s\leq f<t\}$ shrink to the empty set, so finite-measure continuity from above and Proposition \ref{prop:slice-flat-slab-estimate} give $B_s^-\to B_t^-$. If $s\downarrow t$, compare the lower slice at $s$ to the closed-sublevel formula at $t$. Their flat distance is bounded by $\mu(\{t<f<s\})$, which again tends to zero. Thus the right limit is $B_t^+$. The other two limits follow from the same calculation with the endpoint conventions reversed.

Subtract the open- and closed-sublevel formulas to get the jump identity and its bound $\mathcal F_{\mathbb R}(B_t^+-B_t^-)\leq\mu(\{f=t\})$. This gives continuity whenever the latter measure is zero. The finite measure $f_\#\mu$ on $\mathbb R$ has at most countably many atoms: for each positive integer $j$, it has only finitely many atoms of size at least $1/j$, and their union contains every atom. Therefore the possible discontinuity set is countable. A positive level mass is only a sufficient location for a possible jump; the displayed current difference can cancel.
\end{proof}

\begin{proposition}[Bounded variation of the flat-valued slice map]\label{prop:slice-flat-total-variation}
For every finite increasing sequence $t_0<\cdots<t_k$,
\[
 \sum_{i=1}^k\mathcal F_{\mathbb R}(B_{t_i}^- -B_{t_{i-1}}^-)
 \leq\mu(\{t_0\leq f<t_k\})
 \leq\mathbf M(T)+\mathbf M(\partial T).
\]
In particular, for every finite set of levels $F$,
\[
 \sum_{t\in F}\mathcal F_{\mathbb R}(B_t^+-B_t^-)
 \leq\mathbf M(T)+\mathbf M(\partial T).
\]
This is a direct quantitative consequence of the one-sided slicing formulas of \cite{Sim18}, Section 4.
\end{proposition}
\begin{proof}
Apply the slab estimate to each adjacent pair. The half-open slabs are pairwise disjoint and their union is the outer slab, so their measures add exactly. For jumps, use the bound by $\mu(\{f=t\})$ and disjointness of distinct level sets. No completeness assertion for the space of flat currents is needed: all one-sided limits have already been exhibited by explicit currents.
\end{proof}

\begin{corollary}[Flat continuity of almost-everywhere geometric slices]\label{cor:geometric-slice-flat-continuity}
If $T$ is compactly supported and integral, choose the full-measure set of good levels from Proposition \ref{prop:slice-common-exceptional-set} and Proposition \ref{prop:integral-geometric-cut}. At a good level $t$ with $\mu(\{f=t\})=0$, every sequence of good levels $t_j\to t$ satisfies
\[
 \mathcal F_{\mathbb R}(S_{t_j}-S_t)\longrightarrow0.
\]
This conclusion concerns flat convergence, and does not assert convergence of slice masses. It combines the geometric slicing theorem with the exceptional-level distinction in \cite{Sim18}, Lemma 4.5 and equations (4.6)--(4.7).
\end{corollary}
\begin{proof}
At each good level $S_t=B_t^-$. Theorem \ref{thm:slice-flat-one-sided} makes the latter map continuous at the stated non-atomic level. Substituting the equal currents proves the claim, regardless of which side the sequence approaches from.
\end{proof}
